% ==========================================================
\section{Visual Representation Theory}
\label{sec:visual}
% ==========================================================

The choice between textual and graphical representation has measurable consequences for human comprehension and problem-solving efficiency.
Textual representations are inherently one-dimensional and processed serially by the auditory-linguistic system~\cite{larkin_why_1987}.
This requires the reader to hold multiple, potentially distant elements in working memory while searching linearly through the representation.
By contrast, graphical representations are two-dimensional and processed in parallel by the visual system to which approximately a quarter of the brain's total capacity is devoted~\cite{moody_physics_2009}.

The concept of \textit{computational offloading}~\cite{larkin_why_1987} formalises this distinction.
Diagrams index information by spatial location rather than sequential position, allowing elements that must be reasoned about together to be grouped physically.
Spatial, topological, and geometric relationships can thereby be perceived directly, whereas equivalent inferences drawn from text require explicit and effortful computation.
This reduces extraneous cognitive load as identified by Cognitive Load Theory~\cite{sweller_cognitive_2020} and frees working memory for substantive reasoning tasks.

Despite the well-established advantages of graphical representation, the design of visual notations in software engineering has historically been treated as secondary to semantic concerns.
Often it is driven by intuition rather than principled theory~\cite{moody_physics_2009}.
The \textit{Physics of Notations} addresses this deficit by establishing a prescriptive, empirically grounded framework for designing cognitively effective visual languages~\cite{moody_physics_2009}.
The theory derives nine principles from cognitive science and visual perception research~\cite{moody_physics_2009}:
\begin{itemize}
	\item \textbf{Semiotic Clarity:} Each semantic construct should map to one distinct graphical symbol; violations produce redundancy, overload, excess, or deficit.
	\item \textbf{Perceptual Discriminability:} Symbols must be easy to tell apart; discriminability depends on the "visual distance" across variables such as shape, color, and size.
	\item \textbf{Semantic Transparency:} Symbols should visually suggest their meaning so that appearance provides intuitive cues and reduces learning effort.
	\item \textbf{Complexity Management:} Notations should offer mechanisms (for example, modularization, hierarchical grouping, or layered views) to present information incrementally and avoid cognitive overload.
	\item \textbf{Cognitive Integration:} Provide contextual anchors, overview maps, and cross-references so readers can relate individual diagrams to the larger system of representations.
	\item \textbf{Visual Expressiveness:} Leverage multiple visual channels (color, size, brightness, texture, position, etc.) so the notation can encode information across a richer perceptual space.
	\item \textbf{Dual Coding:} Combine concise textual labels or annotations with graphics to leverage both verbal and visual cognitive systems and improve comprehension and retention.
	\item \textbf{Graphic Economy:} Keep the visual vocabulary compact so distinct symbols remain within human discriminative abilities, especially for novice users.
	\item \textbf{Cognitive Fit:} Tailor the notation to task and audience, offering simplified and expert variants so the representation matches users' goals and skills.
\end{itemize}

These principles interact in practice: increasing Visual Expressiveness (e.g., adding color) can improve Perceptual Discriminability and support offset problems introduced by a large graphic vocabulary~\cite{moody_physics_2009}.